\documentclass[11pt,a4paper]{article}

\usepackage{amsmath,amssymb,amsfonts,amsthm}
\usepackage{graphicx}
\usepackage{hyperref}
\usepackage{booktabs}
\usepackage{geometry}
\geometry{margin=2.7cm}

\newtheorem{theorem}{Theorem}
\newtheorem{proposition}{Proposition}
\newtheorem{lemma}{Lemma}

% Shortcuts
\newcommand{\SU}{\mathrm{SU}}
\newcommand{\UU}{\mathrm{U}}
\newcommand{\Nf}{N_f}
\newcommand{\Nc}{N_c}
\newcommand{\Tr}{\mathrm{Tr}}
\newcommand{\adj}{\mathrm{adj}}
\newcommand{\diag}{\mathrm{diag}}
\newcommand{\MeV}{\,\mathrm{MeV}}
\newcommand{\GeV}{\,\mathrm{GeV}}
\newcommand{\TeV}{\,\mathrm{TeV}}
\newcommand{\ie}{i.e.\@}
\newcommand{\eg}{e.g.\@}

\title{A structural sBootstrap EFT from Seiberg confinement:\\
charge--mass duality, Yukawa inversion, and the structural core of the quark mass waterfall}
\author{A.\ Rivero}
\date{\today}

\begin{document}
\maketitle

\begin{abstract}
  We work under the philosophy of the \emph{sBootstrap}~\cite{Rivero:2024epjc}:
  a modern realisation of bootstrap and nuclear-democracy ideas in the spirit
  of Chew~\cite{Chew:1962}, in which the observed fermions are identified with
  confined composite fermions (mesinos and, in a quark--diquark
  complementarity sense, diquarkinos) rather than being endowed with arbitrary
  fundamental
  Yukawa matrices.
  We formulate an exploratory supersymmetric EFT that organises two related
  observations about Standard Model
  flavour data:
  (i) the Koide ratio $K=2/3$ in the charged-lepton sector and the
  approximate inverse Koide relation in the down-quark sector, and
  (ii) the ``waterfall'' pattern whereby consecutive quark-mass triples
  approximately satisfy $K=2/3$ with nontrivial sign assignments.
  The common structural input is a canonically normalised $\UU(3)$ family
  generator $Q_F=z_0 I+T$ whose eigenvalues $q_i$ define a positive operator
  $H_Q=Q_F^\dagger Q_F=\diag(q_1^2,q_2^2,q_3^2)$ with $\Tr(H_Q)=1$.
  For a quadratic charge law $m_i\propto q_i^2$, the Koide ratio becomes a
  trace identity of $H_Q$ and is independent of the Cartan direction; we
  show this direction-independence is unique to the quadratic power.
  We then review $\mathcal{N}{=}1$ $\SU(3)$ SQCD with $\Nf=\Nc=3$:
  the quantum-modified constraint, implemented holomorphically with a
  Lagrange-multiplier superfield and deformed by diagonal mass terms
  $m_i Q^i\tilde Q_i$, yields meson VEVs $\langle M_i\rangle\propto 1/m_i$
  and a $\Lambda$-independent composite spectrum
  $m(M^i{}_j)=m_i m_j/P^{1/3}$ (with $P=m_1m_2m_3$).
  On the mesonic branch, $\det M=\Lambda^6$ implies a holomorphic involution
  $M \leftrightarrow \Lambda^{-2}\adj(M)$ (equivalently $M\to\Lambda^4 M^{-1}$),
  which at the supersymmetric vacuum exchanges the dependence on the
  deformation spurion $m$ and on $m^{-1}$.
  This makes precise how Seiberg mass parameters can generate \emph{inverse}
  Yukawa hierarchies for fields that couple to meson VEVs.
  Finally, we show how the composite doublet $(M_i, \adj(M)_i)$
  has the mathematical structure of an $\mathcal{N}{=}2$ hypermultiplet,
  emergent from confinement plus isospin, and how an effective quartic coupling
  among mesons mixes the direct and inverse branches.  In the
  decoupling limit we prove an interleaving theorem: two monotone branches
  $m=\mu q^2$ and $\hat m=\nu/q^2$ necessarily form an alternating six-rung
  ladder in a calculable window, providing the structural core---rather than a
  quantitative fit---of the quark mass waterfall and of the existence of
  multiple charge decompositions (and sign choices) for the same observed
  masses.  The mesino/lepton map, the quarkino completion, and the
  SUSY-breaking/K\"ahler dressing are treated throughout as EFT assumptions or
  working ans\"atze, not as consequences of SQCD alone; correspondingly, no
  quantitative fit to CKM/PMNS data or proof of full phenomenological viability
  is claimed in this paper.
\end{abstract}

%======================================================================
\section{Introduction}
\label{sec:intro}

The sBootstrap programme~\cite{Rivero:2024epjc} is an attempt to transfer the
bootstrap ethos of strong-interaction physics~\cite{Chew:1962} to the flavour
problem: rather than postulating fundamentally free Yukawa matrices, one seeks
a self-consistent effective description in which the observed spectrum arises
from a small set of symmetry and dynamics inputs.  In the version considered
here, the dynamical core is an $\mathcal{N}{=}1$ confining gauge theory with a
quantum-modified moduli space
whose exact moduli-space relations generate both \emph{direct} and \emph{inverse}
hierarchies from a single set of deformation parameters.
The paper should be read as a controlled \emph{exploratory EFT proposal}:
its goal is to isolate the exact holomorphic algebra, make the gauge-invariant
embedding assumptions explicit, and state clearly which phenomenological steps
are still ans\"atze.

As diagnostics (and as targets for a prospective embedding), we keep track of
three empirical regularities.
First, the Koide relation for charged leptons,
\begin{equation}
  K\bigl(\sqrt{m_e},\sqrt{m_\mu},\sqrt{m_\tau}\bigr)
  \equiv \frac{m_e+m_\mu+m_\tau}{\bigl(\sqrt{m_e}+\sqrt{m_\mu}+\sqrt{m_\tau}\bigr)^2}
  = \frac{2}{3}\,,
  \label{eq:koide_leptons}
\end{equation}
is satisfied to the $10^{-5}$ level with pole masses~\cite{Koide:1983qe,PDG:2024}.
An analogous statement holds approximately for down quarks after inversion
of masses (with a conventional choice of running scheme/scale)~\cite{Rivero:2005}.
Independently, ``waterfall'' relations in which consecutive quark-mass
triples are close to $2/3$ have been observed, but require nontrivial sign
assignments inside the Koide diagnostic~\cite{Zenczykowski:2012}.

This paper develops a single supersymmetric framework that organises these
facts through three ingredients:
\begin{enumerate}
  \item Canonical $\UU(3)$ family normalisation fixes a universal offset
    $z_0=1/\sqrt6$ in the Cartan eigenvalues, producing three charges
    $q_i=z_0+z_i$ with $\sum_i q_i^2=1$ and $\sum_i q_i=\sqrt{3/2}$.
  \item In $\mathcal{N}{=}1$ $\SU(3)$ SQCD with $\Nf=\Nc=3$ (quantum-modified
    moduli space),
    diagonal holomorphic mass deformations $m_i Q^i\tilde Q_i$ generate
    meson VEVs $\langle M_i\rangle\propto 1/m_i$ and a composite spectrum
    in which the leading holomorphic mass scales are simple algebraic functions of the
    \emph{Seiberg masses} $m_i$.
  \item The composite doublet $(M_i, \adj(M)_i)$ exhibits
    $\mathcal{N}{=}2$-like structure, emergent from confinement
    plus isospin.  An effective quartic coupling among mesons
    mixes the direct and inverse branches, producing an interleaved
    mass ladder.
\end{enumerate}

The goal is not to claim a complete UV-to-SM derivation of flavour.
Rather, we isolate which statements are algebraic identities, which are
exact consequences of holomorphy and the quantum-modified constraint, and which
require an assumed embedding of Standard Model fields into the confined sector
(notably through the identification of mesinos with leptons and of
diquarkino-like constituents with quarks, and through non-holomorphic K\"ahler
dressing).
Accordingly, the manuscript should be read as a structural note plus a
gauge-invariant EFT ansatz, not as a completed quantitative flavour model.
More specifically:
\begin{itemize}
  \item Sections~\ref{sec:u3}, \ref{sec:seiberg}, and
    Theorem~\ref{thm:interleaving} contain the hard algebraic/holomorphic
    results.
  \item Sections~\ref{sec:sbootstrap} and~\ref{sec:yukawa} formulate the
    gauge-invariant sBootstrap EFT ansatz and its operator map.
  \item Sections~\ref{sec:pheno} and~\ref{sec:status} state the viability
    conditions and open problems rather than claiming that they are solved.
\end{itemize}

\paragraph{Signed Koide diagnostic.}
For the waterfall discussion it is convenient to define Koide's ratio as
a function of three \emph{real} inputs $x_i$,
\begin{equation}
  K(x_1,x_2,x_3)
  \;\equiv\;
  \frac{x_1^2+x_2^2+x_3^2}{(x_1+x_2+x_3)^2}\,,
  \label{eq:Kdef}
\end{equation}
so that $K(\sqrt{m_1},\sqrt{m_2},\sqrt{m_3})$ reproduces the usual Koide
parameter while relative sign flips correspond to $x_i\to -x_i$.  We also
define the \emph{inverse Koide} diagnostic by applying~$K$ to inverse square
roots,
\begin{equation}
  K^{-1}(m_1,m_2,m_3)
  \;\equiv\;
  K\!\left(\frac{1}{\sqrt{m_1}},\frac{1}{\sqrt{m_2}},\frac{1}{\sqrt{m_3}}\right)
  = \frac{\sum_i m_i^{-1}}{\bigl(\sum_i m_i^{-1/2}\bigr)^2}\,.
  \label{eq:Kinvdef}
\end{equation}
Throughout, the underlying Seiberg masses depend only on $q_i^2$.  Whenever we
discuss relative sign flips in Koide triples, we therefore write them as
explicit factors $s_i=\pm1$ multiplying the signed square roots
$x_i=s_i\sqrt{m_i}$; unless stated otherwise, the sign choices are \emph{not}
reinterpreted as distinct UV charge assignments.

%======================================================================
\section{Canonical \texorpdfstring{$\UU(3)$}{U(3)} charges and the quadratic identity}
\label{sec:u3}

\subsection{Charges from Cartan symmetry breaking}
\label{sec:charges}

Let $\UU(3)$ be a family symmetry with generator
$Q_F=z_0 I + T$ where $T$ lies in the Cartan of $\SU(3)$ and is traceless,
$T=\diag(z_1,z_2,z_3)$ with $\sum_i z_i=0$.
Canonical normalisation fixes
\begin{equation}
  z_0=\frac{1}{\sqrt6}\,,
  \qquad
  \sum_{i=1}^3 z_i^2=\frac{1}{2}\,,
  \label{eq:canonical_norm}
\end{equation}
so that the eigenvalues $q_i$ of~$Q_F$ are
\begin{equation}
  q_i=z_0+z_i\,,
  \qquad
  \sum_i q_i=\sqrt{\frac{3}{2}}\,,
  \qquad
  \sum_i q_i^2=1\,.
  \label{eq:qsum}
\end{equation}
An explicit parametrisation of the Cartan direction is
\begin{equation}
  z_i=\frac{1}{\sqrt3}\cos\!\left(\alpha+\frac{2\pi(i-1)}{3}\right),
  \qquad i=1,2,3,
  \label{eq:zi_alpha}
\end{equation}
with $\alpha$ an angle on the two-dimensional Cartan plane.

The associated positive operator
\begin{equation}
  H_Q\equiv Q_F^\dagger Q_F = Q_F^2 = \diag(q_1^2,q_2^2,q_3^2)
  \label{eq:HQ}
\end{equation}
obeys $\Tr(H_Q)=1$ by~(\ref{eq:qsum}) and will be the universal object
controlling all charge-based mass formulae below.

\subsection{Quadratic mass law and uniqueness}
\label{sec:unique}

Assume a quadratic charge law
\begin{equation}
  m_i=\mu\,q_i^2\,,
  \label{eq:mi_quad}
\end{equation}
with $\mu$ a generation-independent scale.
Then the Koide ratio becomes a trace identity of~$H_Q$:
\begin{equation}
  K\bigl(\sqrt{m_1},\sqrt{m_2},\sqrt{m_3}\bigr)
  = K(q_1,q_2,q_3)
  = \frac{\sum_i q_i^2}{(\sum_i q_i)^2}
  = \frac{1}{3/2}
  = \frac{2}{3}\,,
  \label{eq:K_trace_id}
\end{equation}
independent of the Cartan direction~$\alpha$.

The special role of the quadratic power is not an accident.
Define for integer~$n$ the generalised ratio
\begin{equation}
  K_n(\alpha)
  \;\equiv\;
  \frac{\sum_i (z_0+z_i(\alpha))^n}{\left(\sum_i (z_0+z_i(\alpha))^{n/2}\right)^2}\,,
  \label{eq:Kn_def}
\end{equation}
whenever the numerator and denominator are real.

\begin{theorem}[Uniqueness of the quadratic direction-independence for integer powers]
\label{thm:unique_n2}
  For integer powers $n$, $\UU(3)$ charges $q_i=z_0+z_i$ with
  (\ref{eq:canonical_norm}), and real $K_n(\alpha)$, the function
  $K_n(\alpha)$ is independent of the Cartan direction~$\alpha$
  if and only if $n=2$ (besides the trivial constant $n=0$).
\end{theorem}

\noindent A short proof is given in Appendix~\ref{app:unique}.
For integer $n\ge3$ it follows directly from the presence of the cubic
invariant $\sum_i z_i^3\propto\cos 3\alpha$ in the binomial expansion of
$\sum_i (z_0+z_i)^n$.

%======================================================================
\section{\texorpdfstring{$\mathcal{N}{=}1$}{N=1} confinement with a quantum-modified moduli space}
\label{sec:seiberg}

\subsection{Confined description and vacuum}
\label{sec:confined}

Consider $\mathcal{N}{=}1$ $\SU(3)$ SQCD with $\Nf=\Nc=3$ flavours of
fundamentals $Q^i$ and antifundamentals $\tilde Q_i$ ($i=1,2,3$) and
diagonal mass deformations in the UV superpotential,
\begin{equation}
  W_{\rm UV}\supset \sum_{i=1}^3 m_i\,Q^i\tilde Q_i\,.
  \label{eq:WUV_mass}
\end{equation}
For $\Nf=\Nc$ the theory confines with a quantum-modified moduli space~\cite{Seiberg:1994bz,Csaki:1996zb,IntriligatorSeibergShenker:1995}.
\paragraph{Confinement vs.\ $s$-confinement.}
In standard SQCD terminology, the case $\Nf=\Nc$ has a quantum-modified moduli
space rather than an $s$-confining dynamical superpotential with a smooth
origin.  We implement the holomorphic constraint with a Lagrange-multiplier
superfield $X$ as in~(\ref{eq:WIR}).  In what follows, statements called
``exact'' refer to holomorphic relations implied by the constraint and the
superpotential; physical masses additionally depend on K\"ahler normalisation.
The IR degrees of freedom on the moduli space are gauge-invariant composites:
mesons $M^i{}_j=Q^i\tilde Q_j$ and baryons $B$, $\tilde B$, subject to the
quantum-modified constraint
\begin{equation}
  \det M - B\tilde B = \Lambda^6\,,
  \label{eq:qmc}
\end{equation}
with $\Lambda$ the dynamical scale.
The constraint is enforced by a Lagrange-multiplier chiral field~$X$, giving
\begin{equation}
  W_{\rm IR}
  = X\bigl(\det M - B\tilde B - \Lambda^6\bigr) + \Tr(m\,M)\,,
  \label{eq:WIR}
\end{equation}
On the mesonic branch $B=\tilde B=0$, and in the basis where $m$ is diagonal,
the $F$-term conditions imply that $M$ is diagonal in the vacuum and yield
\begin{equation}
  \langle M_i\rangle = \frac{\Lambda^2\,P^{1/3}}{m_i}\,,
  \qquad
  P\equiv m_1m_2m_3\,,
  \label{eq:Mvev_susy}
\end{equation}
and
\begin{equation}
  X=-\frac{P^{1/3}}{\Lambda^4}\,,
  \label{eq:Xvev}
\end{equation}
which is generation-independent.
Equation~(\ref{eq:Mvev_susy}) is the first, central appearance of
\emph{inversion}: heavy electric masses correspond to small meson VEVs.

\paragraph{Matrix form of the vacuum.}
For an invertible mass matrix $m$ (not necessarily diagonal), the mesonic-branch
$F$-term solution can be written compactly as
\begin{equation}
  \langle M\rangle
  = \Lambda^2\,(\det m)^{1/3}\,m^{-1}\,,
  \qquad
  \langle \adj(M)\rangle
  = \Lambda^4\,(\det m)^{2/3}\,m\,,
  \qquad
  X=-\frac{(\det m)^{1/3}}{\Lambda^4}\,,
  \label{eq:matrix_vacuum}
\end{equation}
which reduces to~(\ref{eq:Mvev_susy})--(\ref{eq:Xvev}) in the diagonal basis
$m=\diag(m_1,m_2,m_3)$.

\subsection{Physical composite spectrum and inversion}
\label{sec:spectrum}

The holomorphic superpotential~(\ref{eq:WIR}) fixes the structure of
fermion and scalar mass terms up to K\"ahler normalisation.
For our purposes, an important special fact is that several physical
mass scales in the confined theory are \emph{$\Lambda$-independent} and
are algebraic functions of the Seiberg masses $m_i$ alone.

To see this most transparently, define $P=m_1m_2m_3$ and consider small
fluctuations of the mesons around the diagonal vacuum.
For off-diagonal mesons $M^i{}_j$ with $i\neq j$, one finds a holomorphic
mass term controlled by the third flavour $k\neq i,j$,
and after canonical normalisation the leading mass scales take the simple
form~\cite{Csaki:1996zb}:
\begin{equation}
  m\bigl(M^i{}_j\bigr)
  = \frac{m_i m_j}{P^{1/3}}
  = \frac{P^{1/3}}{m_k}\,,
  \qquad i\neq j,\ \{i,j,k\}=\{1,2,3\}\,.
  \label{eq:offdiag_mass}
\end{equation}
The absence of explicit $\Lambda$ dependence holds at the holomorphic level;
physical masses are still dressed by K\"ahler normalisation, which is expected
to be flavour-universal at leading order in the mass deformation.
This relation exhibits a \emph{multiple-description} property: the same
leading mass scale can be written either as a product of two Seiberg masses
or as the inverse of the third (up to the common scale~$P^{1/3}$).
The baryon mass scale is similarly simple,
\begin{equation}
  m(B)\sim P^{1/3}\,,
  \label{eq:baryon_mass}
\end{equation}
again independent of~$\Lambda$.
Appendix~\ref{app:seiberg_spectrum} collects a compact derivation.

The quantum constraint~(\ref{eq:qmc}) implies on the mesonic branch
$\det M=\Lambda^6$, which is preserved by the involution
\begin{equation}
  M \longrightarrow \Lambda^4\,M^{-1}\,,
  \label{eq:involution}
\end{equation}
since $\det(\Lambda^4 M^{-1})=\Lambda^{12}/\det M=\Lambda^6$.
The map is holomorphic and well-defined on the locus where $M$ is invertible
(in particular on the mesonic branch); it is not a globally defined symmetry
of the full moduli space including baryonic directions.
Using the adjugate matrix,
\begin{equation}
  \adj(M)\equiv (\det M)\,M^{-1}=\Lambda^6\,M^{-1}\,,
  \label{eq:adj_def}
\end{equation}
the involution can be written holomorphically as $M\leftrightarrow \Lambda^{-2}\adj(M)$.
At the vacuum~(\ref{eq:Mvev_susy}) this exchanges $\langle M_i\rangle\propto 1/m_i$
with $\adj(M)_i\propto m_i$.

\subsection{Charges for each mass and multiple decompositions}
\label{sec:charge_structure}

If the Seiberg masses follow the quadratic charge law~(\ref{eq:mi_quad}),
then the confined spectrum inherits a hierarchy controlled by the same
operator~$H_Q$ but with different \emph{functions} of it.
For instance,
\begin{equation}
  \langle M_i\rangle \propto \frac{1}{m_i}\propto \frac{1}{q_i^2}\,,
  \qquad
  m\bigl(M^i{}_j\bigr)
  = \frac{m_i m_j}{P^{1/3}}
  \propto
  \frac{q_i^2 q_j^2}{(q_1^2 q_2^2 q_3^2)^{1/3}}\,.
  \label{eq:charge_inherited}
\end{equation}
It is useful to introduce dimensionless \emph{effective charges} for these
composites.  We restrict to Cartan directions for which all $q_i>0$ (the regime
relevant to positive mass deformations), so that the geometric mean below is
real and positive.  Writing
\begin{equation}
  \bar q \equiv (q_1 q_2 q_3)^{1/3}\,,
  \label{eq:qbar}
\end{equation}
equation~(\ref{eq:offdiag_mass}) becomes
\begin{equation}
  m\bigl(M^i{}_j\bigr)
  = \mu\left(\frac{q_i q_j}{\bar q}\right)^2
  = \mu\left(\frac{\bar q}{q_k}\right)^2\,.
  \label{eq:eff_charge_mesons}
\end{equation}
Thus each off-diagonal meson mass is determined by an effective charge that
admits \emph{two} natural decompositions:
either as a product $q_i q_j/\bar q$ or as an inverse $\bar q/q_k$.
These are algebraically equivalent (they are related by $P=m_1m_2m_3$) and do
not represent independent charge assignments.
Because the underlying mass law depends only on squares, the same mass is
also compatible with flipping the sign of any participating charge.

The set of functions of $H_Q$ relevant to later sections can be summarised
schematically as:
\begin{center}
{\small
\begin{tabular}{@{}lll@{}}
  \toprule
  Object & Scaling with Seiberg masses & Scaling with family charges \\
  \midrule
  Direct holomorphic masses & $m_i$ & $q_i^2$ \\
  Meson VEVs (mesonic branch) & $1/m_i$ & $1/q_i^2$ \\
  Off-diagonal meson masses & $m_i m_j/P^{1/3}$ & $(q_i q_j/\bar q)^2$ \\
  Baryon mass scale & $P^{1/3}$ & $\bar q^2$ \\
  \bottomrule
\end{tabular}}
\end{center}
The physical interpretation of these entries depends on how Standard Model
fields couple to the confined sector, which we address next.

%======================================================================
\section{The sBootstrap embedding: mesinos as leptons, diquarkinos as quarks}
\label{sec:sbootstrap}

The preceding sections were phrased in standard SQCD language.  The sBootstrap
interpretation~\cite{Rivero:2024epjc} adds a specific EFT philosophy and a
specific identification of (some of) the confined fermions with the observed
Standard Model fermions.

\subsection{Gauge structure and mesino identification}
\label{sec:sbootstrap_mesinos}

The UV gauge group is taken to be
\begin{equation}
  G_{\rm UV} = \SU(3)_F \times \SU(3)_c \times \SU(2)_L \times \UU(1)_Y\,,
  \label{eq:GUV}
\end{equation}
where $\SU(3)_F$ is a family (preon-colour) gauge group that confines at a
scale~$\Lambda_F$, while the SM gauge factors are weakly gauged spectators for
the family confinement dynamics.
In the SQCD formulas of Section~\ref{sec:seiberg} we denote the same dynamical
scale by~$\Lambda\equiv \Lambda_F$.

The $\SU(3)_F$ matter content is organised so that the mesons transform as a
$3\times 3$ matrix under a global $\SU(3)_L\times\SU(3)_R$:
\begin{equation}
  M^\alpha{}_i \sim Q^\alpha \tilde Q_i\,,\qquad
  Q^\alpha \in (\mathbf{3}_F,\mathbf{3}_L),\qquad
  \tilde Q_i \in (\bar{\mathbf{3}}_F,\bar{\mathbf{3}}_R)\,,
  \label{eq:preons}
\end{equation}
so that $M\in(\mathbf{3}_L,\bar{\mathbf{3}}_R)$.

\paragraph{Gauge-invariant UV field content.}
A minimal gauge-invariant field/spurion assignment that reproduces the EFT used
below is
\begin{center}
\small
\begin{tabular}{@{}lllll@{}}
  \toprule
  Field & $\SU(3)_F$ & $\SU(3)_c$ & $\SU(2)_L$ & $\UU(1)_Y$ \\
  \midrule
  $Q^A_a$ ($A=1,2$) & $\mathbf{3}$ & $\mathbf{1}$ & $\mathbf{2}$ & $-1/2$ \\
  $Q^3_a$ & $\mathbf{3}$ & $\mathbf{1}$ & $\mathbf{1}$ & $+1$ \\
  $\tilde Q^a{}_i$ & $\bar{\mathbf{3}}$ & $\mathbf{1}$ & $\mathbf{1}$ & $0$ \\
  $\Sigma_A{}^i$ & $\mathbf{1}$ & $\mathbf{1}$ & $\mathbf{2}$ & $+1/2$ \\
  $\Sigma_3{}^i$ & $\mathbf{1}$ & $\mathbf{1}$ & $\mathbf{1}$ & $-1$ \\
  \bottomrule
\end{tabular}
\end{center}
The fields $Q^\alpha_a,\tilde Q^a{}_i$ are the family-charged preons.
$\Sigma_\alpha{}^i$ is the electroweak-charged deformation spurion (or,
in a dynamical completion, a family-Higgs superfield) transforming as
$(\bar{\mathbf{3}}_L,\mathbf{3}_R)$ under the global
$\SU(3)_L\times\SU(3)_R$.

\paragraph{Index dictionary.}
In Sections~\ref{sec:u3}--\ref{sec:seiberg} we used standard SQCD notation in
which both indices of $M^i{}_j$ label flavour.  In the sBootstrap embedding we
instead interpret the left index as an $\SU(3)_L$ label (containing electroweak
quantum numbers) and the right index as an $\SU(3)_R$ generation label.
The diagonal spurion background in~(\ref{eq:spurionSigma}) selects an
$\SU(3)_L\times\SU(3)_R$-breaking alignment that allows one to work in a basis
where the deformation is effectively diagonal and the component formulas of
Section~\ref{sec:seiberg} apply.
Embedding electroweak symmetry in $\SU(3)_L$ by
\begin{equation}
  \mathbf{3}_L \to \mathbf{2}_{-1/2}\oplus \mathbf{1}_{+1}\,,
  \label{eq:3L_decomp}
\end{equation}
the meson matrix decomposes into electroweak doublets and singlets,
\begin{equation}
  M^\alpha{}_i
  = \underbrace{M^{1,2}{}_i}_{(\mathbf{2}_{-1/2},\,\bar{\mathbf{3}}_R)}
  \;\oplus\;
  \underbrace{M^{3}{}_i}_{(\mathbf{1}_{+1},\,\bar{\mathbf{3}}_R)}\,.
  \label{eq:meson_decomp}
\end{equation}
The central sBootstrap identification is then:
\begin{equation}
  \psi_{M^{1,2}{}_i}\equiv L_i\,,
  \qquad
  \psi_{M^{3}{}_i}\equiv e^c_i\,,
  \label{eq:mesinos_leptons}
\end{equation}
\ie the nine mesino fermions furnish the SM lepton sector (three lepton
doublets plus three charged singlets).  The scalar components of the same
superfields provide Higgs- and slepton-like composites; the details of how one
light Higgs combination is selected are part of the sBootstrap EFT building.

\paragraph{Compact IR dictionary.}
For later operator matching, the relevant confined objects are
\begin{center}
\small
\begin{tabular}{@{}llll@{}}
  \toprule
  SQCD object & $\SU(3)_L\times\SU(3)_R$ & $\SU(2)_L\times\UU(1)_Y$ & sBootstrap role \\
  \midrule
  $M^A{}_i$ & $(\mathbf{3}_L,\bar{\mathbf{3}}_R)$ & $(\mathbf{2},-1/2)$ & $L_i$ mesinos, $H_d$-type scalars \\
  $M^3{}_i$ & $(\mathbf{3}_L,\bar{\mathbf{3}}_R)$ & $(\mathbf{1},+1)$ & $e^c_i$ mesinos, charged scalars \\
  $\adj(M)_A{}^{i}$ & $(\bar{\mathbf{3}}_L,\mathbf{3}_R)$ & $(\mathbf{2},+1/2)$ & $H_u$-type composite operator \\
  $B,\tilde B$ & $(\mathbf{1},\mathbf{1})$ & $(\mathbf{1},0)$ & singlet baryonic superfields \\
  \bottomrule
\end{tabular}
\end{center}
Throughout, there are two layers of normalisation.  The holomorphic Seiberg
operator $M\sim Q\tilde Q$ has the composite/operator dimension appropriate to
the confined description, while the IR chiral superfield used for the mesino
identification is the canonically normalised field obtained after the unknown
K\"ahler metric is specified.  The hierarchy maps written in terms of $M$ and
$\adj(M)$ therefore refer to holomorphic operator dependence; converting them
into physical masses requires the later K\"ahler dressing.

\paragraph{Diagonal mass deformation and alignment.}
In the pure SQCD discussion of Section~\ref{sec:seiberg} we used a diagonal
mass deformation $m_i Q^i\tilde Q_i$.  In the sBootstrap embedding, the two
flavour indices of SQCD correspond to $\SU(3)_L$ (electroweak) and $\SU(3)_R$
(generation).  A diagonal deformation therefore presupposes an
$\SU(3)_L\times\SU(3)_R$-breaking alignment spurion that identifies a preferred
basis in which the deformation is diagonal.  Throughout, we work in that
aligned basis; the resulting residual freedom is the Cartan direction~$\alpha$
entering the charges~(\ref{eq:zi_alpha}).

\paragraph{One meson field and no Yukawa matrices.}
A characteristic sBootstrap simplification is that the generation structure is
carried by the \emph{meson matrix itself} rather than by independent Yukawa
matrices.  Below $\Lambda_F$, a minimal holomorphic EFT can be organised as
\begin{equation}
  W_{\rm sBoot}
  = X\bigl(\det M - B\tilde B - \Lambda_F^6\bigr)
  + \Sigma_\alpha{}^i\,M^\alpha{}_i
  + \frac{y_d}{\Lambda_F}\,\epsilon_{AB}\,\mathcal{O}_d^{A}{}_{i}\,M^{B}{}_{i}
  + \frac{y_u}{\Lambda_F^3}\,\epsilon_{AB}\,\mathcal{O}_u^{A}{}_{i}\,\adj(M)^{B}{}_{i}
  + \cdots\,,
  \label{eq:Wsboot}
\end{equation}
where $y_{u,d}$ are \emph{numbers} rather than matrices,
$\mathcal{O}_d\sim q_L d^c$ and $\mathcal{O}_u\sim q_L u^c$ are
quark-bilinear placeholders, and the family structure sits entirely in the
single meson matrix.  In the aligned background, the VEVs generated by the
deformation obey $\langle M\rangle\propto m^{-1}$ while
$\langle\adj(M)\rangle\propto m$, so the same spurion controls both inverse and
direct hierarchies.

\paragraph{Gauge invariance of the deformation.}
Because $M^\alpha{}_i$ carries electroweak quantum numbers through
(\ref{eq:3L_decomp})--(\ref{eq:meson_decomp}), the diagonal shorthand
$\Tr(mM)$ used in the pure SQCD sections must be understood as the aligned
background of the gauge-invariant coupling
\begin{equation}
  W \supset \Sigma_\alpha{}^i\,M^\alpha{}_i\,,
  \label{eq:spurionSigma}
\end{equation}
with $\Sigma$ transforming in the representation conjugate to~$M$ under the
gauged $\SU(2)_L\times\UU(1)_Y$ subgroup.  As long as $\Sigma$ is treated as a
field/spurion, the Lagrangian is gauge invariant.  Replacing it by a c-number
matrix $m$ is a statement about an \emph{aligned background} in flavour space,
and generically about an electroweak-broken branch: the simple diagonal formulas
of Sections~\ref{sec:seiberg} and~\ref{sec:yukawa} therefore apply after
alignment, not in the fully electroweak-symmetric phase.
The crucial EFT assumption is that $\Sigma$ is the \emph{only} family-breaking
spurion entering the superpotential at leading order.  Once
$\SU(3)_L\times\SU(3)_R$ rotations are used to align its background, the
remaining flavour freedom is reduced to the Cartan direction~$\alpha$ of the
single operator $H_Q$; this is what prevents $\Sigma$ from reintroducing an
arbitrary Yukawa matrix in disguise.

\subsection{Diquarkinos and quarks (complementarity viewpoint)}
\label{sec:sbootstrap_diquarks}

The sBootstrap also aims at a constituent-level identification of the quark
sector with \emph{diquarkino}-like confined fermions, in the spirit of
hadronic supersymmetry (Miyazawa/Catto--G\"ursey)~\cite{Miyazawa:1966,CattoGursey:1985,CattoGursey:1988}
and of the bootstrap/nuclear democracy picture~\cite{Chew:1962}.  A schematic object is a diquark
\begin{equation}
  D \sim \wedge^2 Q\,,
  \label{eq:diquark}
\end{equation}
whose fermionic partner is a diquarkino.  Such objects are not gauge-invariant
operators of the $\Nf=\Nc=3$ quantum-modified theory, but they naturally appear as
effective coloured constituents in a Higgs/complementarity description~\cite{FradkinShenker:1979,BanksRabinovici:1979} in
which baryons are viewed as quark--diquark bound states, $B\sim QD$.

\paragraph{Bare diquarkino obstruction.}
The simplest mesino-as-lepton embedding already fixes the hypercharge gap of a
bare diquarkino.  Writing
$Q\rightarrow \mathbf{2}_u\oplus \mathbf{1}_v$ and
$\tilde Q\rightarrow \mathbf{1}_t$, the mesons decompose as
$M=Q\tilde Q\rightarrow \mathbf{2}_{u+t}\oplus \mathbf{1}_{v+t}$.
Requiring these mesons to furnish $(L_i,e_i^c)$ imposes
$u+t=-1/2$ and $v+t=+1$, hence $v-u=3/2$.
But the bare diquark always transforms as
\begin{equation}
  \wedge^2 Q \rightarrow \mathbf{2}_{u+v}\oplus \mathbf{1}_{2u}
  = \mathbf{2}_{2u+3/2}\oplus \mathbf{1}_{2u}\,,
  \label{eq:diquark_obstruction}
\end{equation}
so the singlet hypercharge is rigidly lower than the doublet one by~$3/2$.
If the doublet is forced to be the SM quark doublet with $Y=+1/6$, the singlet
is unavoidably pushed to $Y=-4/3$ rather than to $+1/3$ or $-2/3$.
This is the clean obstruction to identifying \emph{bare} $\wedge^2 Q$
diquarkinos with the observed quarks.

\paragraph{EPJC completion.}
The EPJC construction~\cite{Rivero:2024epjc} resolves the obstruction not by
insisting on bare diquarks, but by completing the flavour sector into
$\mathbf{15}\oplus\overline{\mathbf{15}}\oplus\mathbf{24}$ of an $\SU(5)$
flavour group.  In the corresponding branchings, the effective coloured sector
contains quark-doublet-like pieces, their conjugates, and an unavoidable exotic
$\pm4/3$ state per generation.  In the present paper we therefore keep
\emph{mesinos = leptons} and \emph{diquarkinos = quarks} only in the
quark--diquark/complementarity sense: the quarkino fields entering the EFT are
effective coloured multiplets of the completion, not exact gauge-invariant
coordinates of the $\Nf=\Nc=3$ moduli space.

\paragraph{Explicit effective operators.}
In the aligned basis, one leading gauge-invariant working ansatz for the
placeholders in~(\ref{eq:Wsboot}) is
\begin{equation}
  W_{q}^{\rm eff}
  = \sum_{i=1}^3\left[
      \frac{y_d}{\Lambda_F}\,\epsilon_{AB}\,
      \mathcal{Q}^{A a}_{i}\,\mathcal{D}^{c}_{a i}\,M^{B}{}_{i}
      +
      \frac{y_u}{\Lambda_F^3}\,\epsilon_{AB}\,
      \mathcal{Q}^{A a}_{i}\,\mathcal{U}^{c}_{a i}\,\adj(M)^{B}{}_{i}
    \right]
    + M_X\,X^{a}_{i}X^{c}_{a i}+\cdots\,,
  \label{eq:Wq_eff}
\end{equation}
where $A,B=1,2$ are $\SU(2)_L$ indices and $a$ is colour.
The effective coloured multiplets
$\mathcal{Q}_i\sim(\mathbf{3},\mathbf{2})_{+1/6}$,
$\mathcal{D}_i^c\sim(\bar{\mathbf{3}},\mathbf{1})_{+1/3}$,
$\mathcal{U}_i^c\sim(\bar{\mathbf{3}},\mathbf{1})_{-2/3}$ are the SM-like
quarkino combinations of the completion, while
$X_i\sim(\mathbf{3},\mathbf{1})_{+4/3}$ and
$X_i^c\sim(\bar{\mathbf{3}},\mathbf{1})_{-4/3}$ denote the exotic pair
required by the obstruction/completion logic.
Equation~(\ref{eq:Wq_eff}) makes the hierarchy map explicit:
$Y_d^{\rm eff}\propto\langle M\rangle\propto m^{-1}$,
$Y_u^{\rm eff}\propto\langle\adj(M)\rangle\propto m$.
The down-sector term is fixed by hypercharge; the up-sector term should be read
as the leading ansatz rather than as a uniquely derived completion, since more
baroque realisations may route the up sector partly through the
$\mathbf{24}$ or through an additional confining sector.
In particular, anomaly closure, the decoupling of the $\pm4/3$ exotics, and the
full consistency of the coloured quarkino sector are delegated to the external
completion and are not claimed as results of the exact $\Nf=\Nc=3$ dynamics.

\begin{center}
\fbox{\parbox{0.93\linewidth}{
\textbf{Assumption / open problem.}
Equation~(\ref{eq:Wq_eff}) is an EFT realisation of the EPJC completion, not a
derivation from the exact gauge-invariant operator basis of the
$\Nf=\Nc=3$ Seiberg theory.  What is universal and derived in this paper is the
two-operator hierarchy map carried by $M$ and $\adj(M)$; deriving the coloured
quarkino multiplets and the heavy $\pm4/3$ exotics from a fully explicit UV
completion remains open.}}
\end{center}

%======================================================================
\section{From Seiberg masses to Yukawa hierarchies}
\label{sec:yukawa}

\subsection{Holomorphic couplings, spurions, and K\"ahler dressing}
\label{sec:holomorphic}

In supersymmetric effective field theory, superpotential mass terms and
Yukawa couplings are both holomorphic parameters.  In particular, the
diagonal mass matrix $m_i$ in~(\ref{eq:WUV_mass}) can be viewed as a
flavour spurion transforming in $(\mathbf{3},\bar{\mathbf{3}})$ under the
UV global $\SU(3)_L\times \SU(3)_R$ symmetry.
The confined dynamics then converts the spurion~$m_i$ into the algebraic
structures~(\ref{eq:Mvev_susy})--(\ref{eq:offdiag_mass}).

However, physical masses in the IR depend on the K\"ahler potential.
Even when holomorphy fixes \emph{which} combinations of $m_i$ appear in the
superpotential, canonical normalisation can introduce (in general,
generation-dependent and off-diagonal) wavefunction factors.
For the present paper, the key point is that holomorphy provides
\emph{hierarchy functions} of $m_i$ that are then dressed by non-holomorphic
effects.  The dressing is essential for mixing (CKM/PMNS), while the
hierarchy functions control the gross ordering patterns such as direct vs
inverse and the interleaving ladder.

\subsection{Two basic Yukawa mechanisms: coupling to \texorpdfstring{$M$}{M} or \texorpdfstring{$\adj(M)$}{adj(M)}}
\label{sec:two_yukawas}

The inversion relation $\langle M_i\rangle\propto 1/m_i$ implies a sharp
statement about how Yukawa hierarchies can arise from Seiberg masses.

\paragraph{sBootstrap form of the hierarchy map.}
In the sBootstrap EFT~(\ref{eq:Wsboot}), the effective Yukawa structures are
controlled by the spurion-induced VEVs,
\begin{equation}
  Y_d^{\rm eff}\;\sim\;\frac{y_d}{\Lambda_F}\,\langle M\rangle \;\propto\; m^{-1}\,,
  \qquad
  Y_u^{\rm eff}\;\sim\;\frac{y_u}{\Lambda_F^3}\,\langle \adj(M)\rangle \;\propto\; m\,,
  \label{eq:Yeff_map}
\end{equation}
up to non-holomorphic K\"ahler dressing (needed, in particular, to generate
mixing).  Using the matrix vacuum relations~(\ref{eq:matrix_vacuum}), this may
be written more explicitly as
\begin{equation}
  Y_d^{\rm eff}\sim
  y_d\,\Lambda_F\,(\det m)^{1/3}\,m^{-1}\,,
  \qquad
  Y_u^{\rm eff}\sim
  y_u\,\Lambda_F\,(\det m)^{2/3}\,m\,.
  \label{eq:Yeff_map_explicit}
\end{equation}
The same aligned spurion therefore generates two hierarchies and no family
matrix is introduced by hand.

\paragraph{Inverse Yukawas from meson VEVs.}
Suppose a chiral field $\psi_i$ couples (in a flavour-diagonal way) to the
meson operator through an effective superpotential term,
\begin{equation}
  W \supset \frac{\lambda}{M_*}\,\psi_i\psi^c_i\,M^i{}_i\,.
  \label{eq:yukawa_M}
\end{equation}
Then, on the mesonic branch, the leading holomorphic mass parameter induced by
the meson VEV is
\begin{equation}
  m_{\psi_i}\sim \frac{\lambda}{M_*}\,\langle M_i\rangle
  \propto \frac{1}{m_i}
  \propto \frac{1}{q_i^2}\,.
  \label{eq:inverse_yukawa}
\end{equation}
In other words, \emph{Yukawa hierarchies can be inverse to Seiberg mass
hierarchies} for fields that couple to $M$.
This is the cleanest sense in which Seiberg mass parameters act as
``Yukawas'' for other sectors: the confined vacuum performs the inversion.

\paragraph{Direct Yukawas from the adjugate.}
Similarly, couplings to the adjugate operator,
\begin{equation}
  W \supset \frac{\hat\lambda}{\hat M_*}\,\psi_i\psi^c_i\,\adj(M)^i{}_i\,,
  \label{eq:yukawa_adj}
\end{equation}
induce masses that scale directly with $m_i$ (up to common factors) because
$\adj(M)_i\propto m_i$ at the vacuum.  Since $\adj(M)$ is holomorphic in $M$
on the constraint surface~(\ref{eq:qmc}), this ``direct'' option is not an
independent structure: it is the other face of the same meson matrix.

\paragraph{Multiple descriptions.}
The presence of additional holomorphic composites such as the off-diagonal
mesons $M^i{}_j$ implies further Yukawa patterns proportional to
$m_i m_j/P^{1/3}$, cf.~(\ref{eq:offdiag_mass}).  Via~(\ref{eq:eff_charge_mesons})
these patterns admit multiple charge decompositions and therefore multiple
interpretations in terms of effective charges.
Here ``multiple decompositions'' does \emph{not} mean multiple independent
charge assignments.  It is the algebraic rewriting forced by
$P=m_1m_2m_3$ and by the constraint surface itself:
\begin{equation}
  \frac{m_i m_j}{P^{1/3}}
  = \mu\left(\frac{q_i q_j}{\bar q}\right)^2
  = \mu\left(\frac{\bar q}{q_k}\right)^2,
  \qquad
  \bar q\equiv(q_1q_2q_3)^{1/3},\qquad \{i,j,k\}=\{1,2,3\}\,.
  \label{eq:multiple_decomposition}
\end{equation}
The same particle can therefore carry a product form $qq'/\bar q$ or an
inverse form $\bar q/q''$ without introducing new family charges.

\subsection{Aligned SUSY breaking as an EFT ansatz for the quadratic triplets}
\label{sec:soft_origin}

The quadratic triplets $m_i\propto (z_0+z_i)^2$ need not be read as detached
numerology, but neither are they derived here from a controlled SUSY-breaking
sector.  What we do instead is exhibit a minimal EFT ansatz in which the same
family operator
$H_Q=Q_F^\dagger Q_F=\diag(q_1^2,q_2^2,q_3^2)$ controls the leading
SUSY-breaking/K\"ahler deformation.  Introducing a SUSY-breaking spurion
$X=\theta^2 F_X$, the aligned correction
\begin{equation}
  \Delta K
  = \frac{X^\dagger X}{M_*^2}\,
    \left[
      c_0\,\Tr(M^\dagger M)
      + c_1\,\Tr(M^\dagger H_Q M)
      + \cdots
    \right]
  \label{eq:DeltaK_HQ}
\end{equation}
induces, in the diagonal basis,
\begin{equation}
  (\Delta m^2_{\rm soft})_i \propto q_i^2 = (z_0+z_i)^2\,.
  \label{eq:soft_qsquare}
\end{equation}
If the dominant holomorphic deformation is aligned with the same family
operator, the simplest effective choice is precisely
$m_i=\mu q_i^2=\mu(z_0+z_i)^2$.
This does not derive the $q_i$ themselves, but it shows how the quadratic
triplets can be made compatible with a single aligned SUSY-breaking spurion in
the low-energy EFT.

The same logic extends to the off-diagonal composite masses.  Once
$m_i\propto q_i^2$, the exact Seiberg relations automatically generate
product and inverse combinations for the \emph{same} physical state through
(\ref{eq:multiple_decomposition}).  In this sense, much of the observed
``multiple $zz'z''$'' structure is induced by the combination of aligned
SUSY breaking and the quantum-modified constraint, rather than by postulating
separate charge sectors for each particle.

From the Volkov--Akulov viewpoint~\cite{VolkovAkulov:1973}, the goldstino
packages the SUSY-breaking order parameter into nonlinear couplings.  In the
present setting, this makes it natural to interpret departures from the pure
holomorphic limit as aligned soft splittings or K\"ahler dressings of a common
family operator, rather than as extra Yukawa matrices.  We stress again that
this is an EFT interpretation, not a derivation of the UV origin of
$m_i\propto(z_0+z_i)^2$.

A suggestive numerical clue is the near-relation
$m_\pi^2-m_\mu^2\approx f_\pi^2$.  In an exactly restored supersymmetric
chiral multiplet the scalar--fermion mass-squared splitting would vanish, so
relations of this kind are naturally read as soft splittings.  We do
\emph{not}, however, identify the physical pion and muon as a literal
unbroken supermultiplet here: the exact $\Nf=\Nc=3$ Seiberg spectrum does not
yet supply the required gauge/global quantum-number assignment.  We use the
relation only as motivation for looking for aligned SUSY-breaking remnants in
the bosonic sector as well as in the fermionic one.

\subsection{Sign choices and the Koide diagnostic}
\label{sec:signs}

Because the fundamental charge law depends on $q_i^2$, each mass $m_i$
is compatible with two signed square roots
\begin{equation}
  x_i=s_i\sqrt{m_i}\,,
  \qquad
  s_i=\pm1\,,
  \label{eq:signed_roots}
\end{equation}
while the underlying charge law remains $m_i=\mu q_i^2$.
When all three entries in a Koide triple come from the same functional form
in~$q$ (\eg three direct entries $\propto q$ or three inverse entries
$\propto 1/q$), the overall sign can be absorbed and no relative sign flips
are needed to obtain $K=2/3$.
Relative signs become meaningful only in \emph{mixed} triples where some
entries scale as $q$ and others as $1/q$; this will be the case in the
waterfall.

%======================================================================
\section{Emergent doublet structure and the quark mass waterfall}
\label{sec:waterfall}

This section is likewise structural.  What is derived below is the existence of
an interleaved direct/inverse ladder once both branches are present; mapping
that ladder quantitatively onto the observed quark windows requires solving the
quartic corrections and the non-holomorphic dressing, which we do not attempt
here.

\subsection{The effective quartic from the composite doublet}
\label{sec:N2_descent}

The $\Nf=\Nc$ confined theory contains both $m_i$ and $1/m_i$
structures but, by itself, does not provide a controlled coupling between
them.  The composite spectrum, however, has a natural doublet structure:
for each family~$i$, the pair $(M_i, \adj(M)_i)$ furnishes two mass
eigenvalues with opposite charge dependence.  This pairing is electroweak
isospin and has the mathematical form of an $\mathcal{N}{=}2$
hypermultiplet~\cite{APS:1997,SeibergWitten:1994a}---emergent from
confinement plus isospin, not a UV input.

The coupling between branches can be parametrised by introducing an
auxiliary adjoint chiral field~$\Phi$ with mass~$\mu_\Phi$ and the
standard coupling $W \supset \sqrt2\,g\,\tilde Q\,\Phi\,Q$, then
integrating $\Phi$ out via its $F$-term.
One obtains an effective $\mathcal{N}{=}1$ superpotential with a quartic remnant,
\begin{equation}
  W_{\rm eff}
  = -\frac{g^2}{\mu_\Phi}\left[\Tr(Q\tilde Q)^2-\frac{1}{3}\bigl(\Tr Q\tilde Q\bigr)^2\right]
  + \sum_i m_i Q^i\tilde Q_i\,,
  \label{eq:Weff}
\end{equation}
which vanishes as $\mu_\Phi\to\infty$ (decoupling).
Here $\Phi$ is a bookkeeping device for the quartic, not a physical UV
field; the scale~$\mu_\Phi$ parametrises the strength of the inter-branch
coupling in the composite dynamics.

\subsection{Residual quartic in the confined phase}
\label{sec:quartic}

In the confined description $Q\tilde Q\to M$, the quartic becomes
\begin{equation}
  W_{\rm quartic}
  = -\frac{g^2}{\mu_\Phi}\left[\Tr M^2-\frac{1}{3}(\Tr M)^2\right].
  \label{eq:quartic}
\end{equation}
This term couples the diagonal mesons and therefore perturbs the pure
inversion solution~(\ref{eq:Mvev_susy}).  Since $\langle M_i\rangle\propto 1/m_i$,
the fractional correction induced by~(\ref{eq:quartic}) is parametrically
largest for the lightest family and smallest for the heaviest.  This
hierarchy of corrections will reappear as a hierarchy of deviations from
$K=2/3$ in mixed waterfall windows.

\subsection{Interleaving theorem}
\label{sec:interleaving}

In the decoupling limit, the two basic hierarchies are the direct and inverse
branches
\begin{equation}
  a_i=\mu\,q_i^2\,,
  \qquad
  d_i=\nu\,q_i^{-2}\,,
  \label{eq:branches}
\end{equation}
with opposite orderings for $q_1<q_2<q_3$:
$a_1<a_2<a_3$ while $d_1>d_2>d_3$.

\begin{theorem}[Interleaving]
\label{thm:interleaving}
  The six unmixed masses $\{a_1,a_2,a_3,d_1,d_2,d_3\}$ form the alternating ladder
  \begin{equation}
    d_1 > a_3 > d_2 > a_2 > d_3 > a_1
    \label{eq:ladder}
  \end{equation}
  if and only if
  \begin{equation}
    \max\!\bigl(q_1^2 q_3^2,\; q_2^4\bigr)
    \;<\; \frac{\nu}{\mu}
    \;<\; q_2^2\,q_3^2\,.
    \label{eq:window}
  \end{equation}
\end{theorem}

\begin{proof}
  The adjacent inequalities $d_1>a_3>a_2>\cdots>a_1$ reduce to the three
  independent conditions
  $\nu/\mu>q_1^2q_3^2$, $\nu/\mu<q_2^2q_3^2$, and $\nu/\mu>q_2^4$.
  The remaining inequalities follow from $q_1<q_2<q_3$.
\end{proof}

The ladder alternates strictly between inverse and direct entries:
$(d_1,a_3,d_2,a_2,d_3,a_1)$.  Each consecutive triple therefore mixes
functional dependences $q$ and $1/q$ at the level of signed square roots,
and the Koide diagnostic becomes sensitive to relative signs.

\subsection{Waterfall windows and the sign flip}
\label{sec:waterfall_windows}

In the Koide diagnostic, the two branches correspond to
\begin{equation}
  \sqrt{a_i}=\sqrt{\mu}\,q_i\,,
  \qquad
  \sqrt{d_i}=\frac{\sqrt{\nu}}{q_i}\,.
  \label{eq:sqrt_branches}
\end{equation}
For a pure-branch triple, all three entries scale as the same function of $q$
and $K=2/3$ reduces to the trace identity~(\ref{eq:K_trace_id}).
For a mixed triple (\eg $(a_3,d_2,a_2)$), the sum in the denominator of~(\ref{eq:Kdef})
contains both $q$ and $1/q$ contributions, and a relative sign is generically
needed for the ratio to be close to $2/3$.
Empirically, the quark-mass waterfall uses a sign flip in the $(b,c,s)$ window
of physical masses~\cite{Zenczykowski:2012}.  In the present framework, such a
sign flip is interpreted only as a qualitative diagnostic of the residual
quartic mixing
(\ref{eq:quartic}), which can shift the intermediate eigenvalues away from the
unmixed $a_i,d_i$ and thereby modify which branch character dominates each
physical mass.  We do not claim here that the quartic term has already been
shown quantitatively to reproduce the observed windows or sign assignments.
Section~\ref{sec:waterfall} should therefore be read as establishing a
structural motif for the waterfall, not yet a demonstrated quark-mass
mechanism.

%======================================================================
\section{Phenomenology and EFT consistency}
\label{sec:pheno}

The purpose of the present paper is to formulate a controlled flavour EFT, not
to claim that all phenomenology is already computed.  The minimal consistency
checklist is therefore as follows.

\paragraph{Multi-Higgs structure and FCNC.}
The confined meson sector contains three $H_d$-type doublets $M^{A}{}_i$ and,
through $\adj(M)$, three $H_u$-type doublets.  With family-dependent VEVs or
Yukawa textures, generic neutral-scalar exchange induces tree-level FCNC and
new CP phases.  A viable low-energy EFT therefore requires decoupling,
alignment, or natural flavour conservation; at most one may presently invoke a
Cheng--Sher-type suppression as a \emph{working assumption}, not as a derived
result~\cite{Branco:2011iw,ChengSher:1987}.

\paragraph{What the scale $\Lambda_F$ does and does not mean.}
The scale $\Lambda_F$ is the hidden family-confinement scale entering the
holomorphic Seiberg relations.  It should not automatically be identified with
the direct contact-interaction scale probed by compositeness searches on the
observed quarks and leptons.  If the quark sector is partially composite or
quarkino-like, the relevant phenomenological bound applies to the canonically
normalised resonance masses and mixing operators, which must lie well above the
weak scale even if the holomorphic parameter entering the moduli-space algebra
is lower~\cite{Kaplan:1991dc,Strassler:1995,NelsonStrassler:2000}.

\paragraph{Exotics.}
The quark completion that removes the bare-diquarkino obstruction introduces an
exotic $\pm4/3$ coloured state per generation.  In the EFT used here, these
states are assumed heavy (or effectively vector-like after completion) and only
their SM-like mixtures enter the leading operators~(\ref{eq:Wq_eff}).  Any
fully explicit completion must state their masses, anomaly role, and mixing
pattern; they are evidence for extra structure, not decorative bystanders.

\paragraph{Mixing matrices.}
Because the holomorphic structures are functions of the same aligned spurion
$m$, they commute and do not generate CKM or PMNS mixing by themselves.
Mixing must therefore arise from non-holomorphic K\"ahler terms, soft
misalignment, or elementary--composite mixing matrices in the quarkino sector.
This is the honest status of the CKM formulas: until the relevant K\"ahler or
mixing operators are derived, they remain empirical guides rather than EFT
predictions.  Accordingly, the present paper does not claim a derivation of the
CKM or PMNS matrices.  The neutrino sector is even less developed; large PMNS
angles would have to come from baryonic/K\"ahler threshold effects or else be
treated as outside the present scope.

\paragraph{Bosonic tuples.}
The present paper mostly tracks the fermionic tuples, but the same Koide-type
regularities also appear empirically in mesonic triples such as
$(0,\pi,D)$ and $(\pi,D,B)$.  We take that bosonic parallel as a clue that the
same aligned SUSY-breaking/compositeness logic may extend beyond the fermion
sector, but we do not yet derive those meson tuples inside the EFT.
Absent the missing K\"ahler/completion analysis, the paper should therefore be
read as a structural flavour-EFT proposal rather than as a viability proof.

%======================================================================
\section{Status, interpretation, and open problems}
\label{sec:status}

The results above split naturally into five classes:
\begin{enumerate}
  \item \textbf{Algebraic identities:} the trace identity~(\ref{eq:K_trace_id})
    and the uniqueness Theorem~\ref{thm:unique_n2} depend only on canonical
    $\UU(3)$ normalisation and on the assumed monomial charge law.
  \item \textbf{Exact holomorphic constraint relations:}
    the mesonic-branch vacuum $\langle M\rangle\propto m^{-1}$ and the composite
    spectrum relations (\ref{eq:offdiag_mass})--(\ref{eq:baryon_mass}) follow
    from the constraint $\det M-B\tilde B=\Lambda^6$ and its holomorphic
    implementation in~(\ref{eq:WIR}).  Physical masses additionally depend on
    K\"ahler normalisation.
  \item \textbf{sBootstrap postulates:} the identification of mesinos with
    leptons~(\ref{eq:mesinos_leptons}), the quark--diquark/complementarity
    interpretation of quarks as diquarkinos (Section~\ref{sec:sbootstrap_diquarks}),
    and the existence of an alignment spurion selecting a diagonal deformation
    basis are assumptions of the sBootstrap EFT rather than consequences of SQCD.
  \item \textbf{Model-dependent embeddings and mixing:} mapping the holomorphic
    structures to SM pole masses (including CKM/PMNS) requires specifying the
    quark-sector operator content in~(\ref{eq:Wq_eff}) and controlling the
    non-holomorphic K\"ahler dressing.  This is where mixing matrices arise,
    since functions of a single diagonal operator $H_Q$ commute and cannot
    generate CKM without non-holomorphic effects.
  \item \textbf{Speculative UV motivation:} the reading of
    $m_i\propto(z_0+z_i)^2$ as an aligned SUSY-breaking/K\"ahler pattern
    (Section~\ref{sec:soft_origin}) is a concrete EFT interpretation, while
    deeper origins in charge-lattice/BPS or M-theory constructions are
    conjectural and are not used as premises in the derivations above
    \cite{Narain:1986,Ferrara:1995ih}.
\end{enumerate}

Earlier versions of this work were presented as two short notes
(``Quadratic charge--mass relation from gauged family symmetry'' and
``The quark mass waterfall from emergent interleaving''),
which separated the quadratic charge identity and the interleaving theorem.
In this consolidated paper we have focused instead on the structural and
supersymmetric origin of: (i) the direct/inverse pairing, (ii) the existence of
multiple charge decompositions for the same composite masses, and (iii) the
interleaving theorem that forces a mixed ordering whenever both branches are
present and $\nu/\mu$ lies in the window~(\ref{eq:window}).

Important open problems include:
(a) computing (or parametrising with control) the relevant K\"ahler metrics and
their flavour dependence;
(b) deriving the coloured quarkino completion and the $\pm4/3$ exotics from an
explicit UV model rather than from EFT placeholders;
(c) understanding whether the quartic remnant~(\ref{eq:quartic}) can generate the
observed pattern of sign choices in the waterfall windows;
(d) clarifying whether the same aligned SUSY-breaking logic can account for the
empirical meson tuples $(0,\pi,D)$ and $(\pi,D,B)$ as well as for the fermionic
ones.

%======================================================================
\section{Conclusions}
\label{sec:concl}

We have outlined a structural SUSY EFT proposal in which canonical $\UU(3)$
family normalisation, $\mathcal{N}{=}1$ confinement with a quantum-modified
moduli space, and the emergent composite doublet structure
jointly organise why:
(i) a quadratic charge law yields an exact $K=2/3$ identity independent of
the Cartan direction,
(ii) Seiberg mass deformations generate \emph{inverse} hierarchies through
meson VEVs and a $\Lambda$-independent composite spectrum with multiple
charge decompositions, and
(iii) the coexistence of direct and inverse branches forces an interleaved
mass ladder in a calculable parameter window, providing the structural core
of the quark mass waterfall.
In the EFT reading advocated here, the same family operator $H_Q$ can also be
used to organise the leading aligned SUSY-breaking/K\"ahler deformations, so
that the triplets $(z_0+z_i)^2$ are not merely postulated but are compatible
with the minimal quadratic symmetry-breaking pattern.  Embedding these structures into
the full SM flavour sector remains model-dependent, primarily controlled by
K\"ahler dressing, by the explicit quarkino/exotic completion, and by whether
the bosonic tuples can be brought under the same umbrella as the fermionic
ones.  No claim is made here that the present EFT already supplies a
phenomenologically complete or uniquely UV-derived flavour model.

%======================================================================
\appendix

\section{Proof of Theorem~\ref{thm:unique_n2}}
\label{app:unique}

With $q_i=z_0+z_i$ and $z_i$ parametrised by~(\ref{eq:zi_alpha}),
one has $\sum_i z_i=0$ and $\sum_i z_i^2=1/2$ by construction, while
\begin{equation}
  \sum_i z_i^3 \;=\; \frac{1}{3\sqrt3}\sum_i \cos^3\!\left(\alpha+\frac{2\pi(i-1)}{3}\right)
  \;\propto\; \cos 3\alpha\,,
  \label{eq:sumz3}
\end{equation}
so the cubic power sum depends on the Cartan direction.
For integer $n\ge3$, expanding $\sum_i q_i^n=\sum_i(z_0+z_i)^n$ by the
binomial theorem, the term $\binom{n}{3}z_0^{n-3}\sum_i z_i^3$ appears with
nonzero coefficient and introduces $\alpha$ dependence into the numerator of
$K_n(\alpha)$ in~(\ref{eq:Kn_def}).  For $n=2$ one has $\sum_i q_i^2=3z_0^2+\sum_i z_i^2=1$
and $(\sum_i q_i)^2=(3z_0)^2=3/2$, giving $K_2=2/3$ independent of~$\alpha$.
For $n=1$ the denominator contains $\sum_i q_i^{1/2}$ and is $\alpha$ dependent.
The case $n=0$ gives the trivial constant $K_0=1/3$.
For non-integer real $n$, the same cubic-power-sum mechanism suggests that
$\alpha$ dependence should persist generically, but we do not claim a theorem:
the non-integer extension is heuristic and is not used elsewhere in the paper.

\section{Derivation of the quantum-modified spectrum relations}
\label{app:seiberg_spectrum}

We sketch the derivation of~(\ref{eq:offdiag_mass}) following the standard
quantum-modified effective description~(\ref{eq:WIR}).
On the diagonal vacuum $M=\diag(M_1,M_2,M_3)$ with $B=\tilde B=0$, the $F$-term
conditions imply $M_i=\Lambda^2 P^{1/3}/m_i$ and $X=-P^{1/3}/\Lambda^4$.

For an off-diagonal meson $M^i{}_j$ with $i\neq j$, the relevant holomorphic
mass term comes from the second derivative of the determinant,
$\partial^2(\det M)/\partial M^i{}_j \partial M^j{}_i = M_k$, where $k$ is the
third index.  Canonical normalisation for mesons is $M/\Lambda$ at leading order,
so the leading mass scale is
\begin{equation}
  m\bigl(M^i{}_j\bigr)
  \sim \Lambda^2\,|X M_k|
  = \Lambda^2\left|\frac{P^{1/3}}{\Lambda^4}\frac{\Lambda^2 P^{1/3}}{m_k}\right|
  = \frac{P^{2/3}}{m_k}
  = \frac{m_i m_j}{P^{1/3}}\,,
\end{equation}
using $P=m_1m_2m_3$.  The overall proportionality depends on the K\"ahler metric
but is generation-independent at leading order; the dependence on $m_i$ is fixed
by holomorphy.

\begin{thebibliography}{99}

\bibitem{Koide:1983qe}
  Y.~Koide,
  Phys.\ Lett.\ B \textbf{120}, 161 (1983).

\bibitem{Rivero:2024epjc}
  A.~Rivero,
  Eur.\ Phys.\ J.\ C \textbf{84}, 1058 (2024),
  \href{https://doi.org/10.1140/epjc/s10052-024-13368-3}%
    {\texttt{doi:10.1140/epjc/s10052-024-13368-3}}.

\bibitem{Chew:1962}
  G.F.~Chew,
  Rev.\ Mod.\ Phys.\ \textbf{34}, 394 (1962),
  \href{https://doi.org/10.1103/RevModPhys.34.394}%
    {\texttt{doi:10.1103/RevModPhys.34.394}}.

\bibitem{ChewFrautschi:1962}
  G.F.~Chew and S.C.~Frautschi,
  Phys.\ Rev.\ Lett.\ \textbf{8}, 41 (1962),
  \href{https://doi.org/10.1103/PhysRevLett.8.41}%
    {\texttt{doi:10.1103/PhysRevLett.8.41}}.

\bibitem{Miyazawa:1966}
  H.~Miyazawa,
  Prog.\ Theor.\ Phys.\ \textbf{36}, 1266 (1966),
  \href{https://doi.org/10.1143/PTP.36.1266}%
    {\texttt{doi:10.1143/PTP.36.1266}}.

\bibitem{CattoGursey:1985}
  S.~Catto and F.~G\"ursey,
  Nuovo Cim.\ A \textbf{86}, 201 (1985),
  \href{https://doi.org/10.1007/BF02902548}%
    {\texttt{doi:10.1007/BF02902548}}.

\bibitem{CattoGursey:1988}
  S.~Catto and F.~G\"ursey,
  Nuovo Cim.\ A \textbf{99}, 685 (1988),
  \href{https://doi.org/10.1007/BF02730633}%
    {\texttt{doi:10.1007/BF02730633}}.

\bibitem{FradkinShenker:1979}
  E.~Fradkin and S.H.~Shenker,
  Phys.\ Rev.\ D \textbf{19}, 3682 (1979),
  \href{https://doi.org/10.1103/PhysRevD.19.3682}%
    {\texttt{doi:10.1103/PhysRevD.19.3682}}.

\bibitem{BanksRabinovici:1979}
  T.~Banks and E.~Rabinovici,
  Nucl.\ Phys.\ B \textbf{160}, 349 (1979).

\bibitem{PDG:2024}
  S.~Navas \textit{et al.} (Particle Data Group),
  Phys.\ Rev.\ D \textbf{110}, 030001 (2024).

\bibitem{Rivero:2005}
  A.~Rivero,
  arXiv:hep-ph/0505220 (2005).

\bibitem{Zenczykowski:2012}
  P.~Zenczykowski,
  Phys.\ Rev.\ D \textbf{86}, 117303 (2012).

\bibitem{Seiberg:1994bz}
  N.~Seiberg,
  Nucl.\ Phys.\ B \textbf{435}, 129 (1995).

\bibitem{Csaki:1996zb}
  C.~Cs\'aki, M.~Schmaltz, and W.~Skiba,
  Phys.\ Rev.\ Lett.\ \textbf{78}, 799 (1997).

\bibitem{IntriligatorSeibergShenker:1995}
  K.~Intriligator, N.~Seiberg, and S.H.~Shenker,
  Phys.\ Lett.\ B \textbf{342}, 152 (1995).

\bibitem{SeibergWitten:1994a}
  N.~Seiberg and E.~Witten,
  Nucl.\ Phys.\ B \textbf{426}, 19 (1994);
  \textbf{431}, 484 (1994).

\bibitem{APS:1997}
  P.C.~Argyres, M.R.~Plesser, and A.D.~Shapere,
  Nucl.\ Phys.\ B \textbf{483}, 172 (1997),
  arXiv:hep-th/9705069.

\bibitem{VolkovAkulov:1973}
  D.V.~Volkov and V.P.~Akulov,
  Phys.\ Lett.\ B \textbf{46}, 109 (1973).

\bibitem{Narain:1986}
  K.S.~Narain,
  Phys.\ Lett.\ B \textbf{169}, 41 (1986).

\bibitem{Ferrara:1995ih}
  S.~Ferrara, R.~Kallosh, and A.~Strominger,
  Phys.\ Rev.\ D \textbf{52}, 5412 (1995),
  arXiv:hep-th/9508072.

\bibitem{Kaplan:1991dc}
  D.B.~Kaplan,
  Nucl.\ Phys.\ B \textbf{365}, 259 (1991).

\bibitem{Strassler:1995}
  M.J.~Strassler,
  Phys.\ Lett.\ B \textbf{376}, 119 (1996),
  arXiv:hep-ph/9510342.

\bibitem{NelsonStrassler:2000}
  A.E.~Nelson and M.J.~Strassler,
  JHEP \textbf{09}, 030 (2000),
  arXiv:hep-ph/0006251.

\bibitem{Branco:2011iw}
  G.C.~Branco, P.M.~Ferreira, L.~Lavoura, M.N.~Rebelo,
  M.~Sher, and J.P.~Silva,
  Phys.\ Rept.\ \textbf{516}, 1 (2012),
  arXiv:1106.0034.

\bibitem{ChengSher:1987}
  T.P.~Cheng and M.~Sher,
  Phys.\ Rev.\ D \textbf{35}, 3484 (1987).

\end{thebibliography}

\clearpage
\tableofcontents

\end{document}
